\documentclass[11pt]{article}

\usepackage[a4paper,margin=30mm]{geometry}
\usepackage[T1]{fontenc}
\usepackage{amsmath}
\usepackage{xcolor}
\usepackage{hyperref}

\definecolor{paperblue}{HTML}{245CA6}
\hypersetup{
  colorlinks=true,
  linkcolor=paperblue,
  urlcolor=paperblue,
  pdfborder={0 0 0},
  pdftitle={Lean Paper Project Template}
}

% The installed VS Code extension intercepts this private URI locally.
% No web request or operating-system protocol registration is involved.
\newcommand{\leanlink}[2]{%
  \href{lean-paper:#1}{#2}%
}

% Replace this simple heading helper with the theorem style used by your paper.
% Keeping the link on a short heading avoids page-break ambiguity.
\newenvironment{leanstatement}[2]{%
  \par\medskip\noindent
  \leanlink{#1}{\textbf{#2}}%
  \quad\itshape
}{%
  \par\medskip
}

\title{Lean Paper Project Template}
\author{Replace with author name}
\date{}

\begin{document}
\maketitle

\section{Linked statements}

The blue headings below are semantic PDF links. After opening this PDF with
Lean Paper Viewer, clicking a heading opens the corresponding declaration in
the right-hand Lean editor.

\begin{leanstatement}{def:twice}{Definition 1.1 (Twice).}
For a natural number $n$, define
\[
  \operatorname{twice}(n) := n+n.
\]
\end{leanstatement}

\begin{leanstatement}{thm:twice-eq-add-self}{Theorem 1.2 (Expansion of twice).}
For every natural number $n$,
\[
  \operatorname{twice}(n)=n+n.
\]
\end{leanstatement}

\begin{leanstatement}{thm:zero-add}{Theorem 1.3 (Zero on the left).}
For every natural number $n$,
\[
  0+n=n.
\]
\end{leanstatement}

\section{Links in ordinary prose}

Links need not be attached to theorem headings. For example, the term
\leanlink{def:twice}{\emph{twice}} refers to the function defined above.
In an ordinary sentence, we can also link the assertion that
\leanlink{thm:zero-add}{adding zero on the left leaves a natural number unchanged}.
A short inline formula can be linked as well:
\leanlink{thm:twice-eq-add-self}{\(\operatorname{twice}(n)=n+n\)}.

These links reuse the same semantic IDs as the headings, so they open the
same Lean declarations without additional manifest entries.

\section{Adapting the template}

Keep \verb|\leanlink| in the preamble, but adapt \verb|leanstatement| to the
theorem style of the actual paper. Use stable semantic IDs and update
\texttt{lean-paper-viewer.json} whenever a link is added or removed.

\end{document}
